%!TEX root main.tex

\section{Polynomial vs.~Hadamard automata}
\label{app:polynomial automata}
% lax begin Lax946791.PolynomialAutomata

In this section, we give more details about the relationship
between polynomial and Hadamard automata,
expanding on~\cref{rem:Hadamard and polynomial automata}.
%
% Our presentation of polynoimial automata follows~\cite{Bojanczyk:AutomataToolbox:2025}.
%
A \emph{polynomial automaton}~\cite{BenediktDuffSharadWorrell:PolyAut:LICS:2017} is a tuple
$\AA = \tuple{k, \Sigma, Q, q_I, \Delta, F}$
where $k \in \N$ is the \emph{dimension},
$\Sigma$ is a finite \emph{input alphabet},
$Q = \Q^k$ is the set of \emph{states},
$q_I \in Q$ is the \emph{initial state},
$\Delta : \Sigma \to Q \to Q$ is the polynomial \emph{transition function},
and $F : Q \to \Q$ is the polynomial \emph{output function}.
%
For every input symbol $a \in \Sigma$,
the polynomial map $\Delta^a : Q^k \to Q^k$
is presented as a tuple of polynomials $\Delta^a = \tuple{\Delta^a_1, \dots, \Delta^a_k} \in \poly \Q {y_1, \dots, y_k}^k$,
inducing a polynomial action on states
%
\begin{align*}
    q \cdot a := \Delta^a(q) = \tuple{\Delta^a_1(q), \dots, \Delta^a_k(q)} \in Q,
    \quad\text{for all } q \in Q.
\end{align*}
%
Similarly, the polynomial output function $F$ is presented as a polynomial $F \in \poly \Q {y_1, \dots, y_k}$
inducing the corresponding polynomial map
$F(q) = F(q_1, \dots, q_k) \in \Q$ for every $q = \tuple{q_1, \dots, q_k} \in Q$.
%
The action of $\Sigma$ is extended to words $w \in \Sigma^*$ homomorphically:
$q \cdot \e := q$ and $q \cdot (a \cdot w) := (q \cdot a) \cdot w$.
%
The \emph{semantics} of a state $q \in Q$ of a polynomial automaton $\AA$
is the series defined as follows:
%
\begin{align*}
    \Asem \AA q &\in \series \Q \Sigma \\
    \coefficient w (\Asem \AA q) &:= F(q \cdot w), \quad \text{for every } w \in \Sigma^*.
\end{align*}
%
The semantics of the automaton $\AA$ is the series recognised by the initial state
$\Asem \AA {q_I} {}$.
%
The following lemma shows that polynomial automata recognise precisely the Hadamard-finite series,
up to reversal of the input.
%
Since Hadamard-finite series coincide with the series recognised by Hadamard automata
(\cref{lem:Hadamard automata and series}),
the same holds for Hadamard automata.
%
\begin{lemma}
    A series $f \in \series \Q \Sigma$ is recognisable by a polynomial automaton
    if, and only if, its reversal $f^R$ is Hadamard finite
    (equivalently, recognised by a Hadamard automaton).
\end{lemma}
%
\noindent
Recall that $f^R$ is the series that maps a word $w \in \Sigma^*$ to $f(w^R)$,
where $w^R = a_n \cdots a_1$ is the reversal of $w = a_1 \cdots a_n$.
%
\begin{proof}
    For the ``only if'' direction,
    let $\AA$ be a polynomial automaton recognising
    the series $f = \Asem \AA {q_I}$.
    %
    For every $1 \leq i \leq k$,
    consider the series $g_i \in \series \Q \Sigma$ defined as
    %
    \begin{align*}
        g_i(w) := \pi_i(q_I \cdot w^R),
        \quad \text{for all } w \in \Sigma^*.
    \end{align*}
    %
    In other words, $g_i(w)$ is the value of component $i$
    after reading $w^R$ from the initial state.
    %
    Consider now the Hadamard algebra
    %
    \begin{align*}
        A := \poly \Q {g_1, \dots, g_k}_{\hadamard} \subseteq \series \Q \Sigma.
    \end{align*}
    %
    Thanks to the working characterisation of Hadamard-finite series from~\cref{lem:Hadamard finite - working definition},
    this part of the proof is concluded by the following two claims.
    %
    \begin{claim*}
        $f^R$ is in $A$.
    \end{claim*}
    \begin{proof}
        By definition of $f$, we have $f = \Asem \AA {q_I} {}$.
        Consequently for every $w \in \Sigma^*$ we have
        %
        \begin{align*}
            f^R(w)
                &= F(q_I \cdot w^R)
                = F(\pi_1(q_I \cdot w^R), \dots, \pi_k(q_I \cdot w^R))
                = F(g_1(w), \dots, g_k(w)),
        \end{align*}
        and thus $f^R = F(g_1, \dots, g_k)$ is in $A$,
        as required.
    \end{proof}

    \begin{claim*}
        For every $a \in \Sigma$ and $1 \leq i \leq k$,
        we have $\deriveleft a g_i \in A$.
    \end{claim*}
    \begin{proof}
        By the definition of $g_i$, for every $w \in \Sigma^*$ we have
        %
        \begin{align*}
            (\deriveleft a g_i)(w)
                &= g_i(a \cdot w)
                = \pi_i(q_I \cdot (a \cdot w)^R)
                = \pi_i(q_I \cdot (w^R \cdot a))
                = \pi_i((q_I \cdot w^R) \cdot a)
                = \Delta^a_i(q_I \cdot w^R)
                = \Delta^a_i(g_1(w), \dots, g_k(w)),
        \end{align*}
        and thus $\deriveleft a g_i = \Delta^a_i(g_1, \dots, g_k)$
        is in $A$, as required.
    \end{proof}

    For the ``if'' direction, let $g \in \series \Q \Sigma$
    belong to a finitely generated Hadamard algebra $A := \poly \Q {g_1, \dots, g_k}_{\hadamard}$
    closed under left derivatives $\deriveleft a$ (with $a \in \Sigma$).
    %
    Since $g \in A$, we can write $g = F(g_1, \dots, g_k)$
    for some polynomial $F \in \poly \Q {y_1, \dots, y_k}$.
    %
    For every $1 \leq i \leq k$ and $a \in \Sigma$,
    the series $\deriveleft a g_i$ is in $A$ and thus it can be written as
    $\Delta^a_i(g_1, \dots, g_k)$ for some polynomial $\Delta^a_i \in \poly \Q {y_1, \dots, y_k}$.
    %
    Finally, let $q_I := \tuple{g_1(\e), \dots, g_k(\e)} \in \Q^k$ be the tuple of constant terms of the generators.
    %
    This provides us with the data required to construct a dimension $k$ polynomial automaton
    $\AA = \tuple{k, \Sigma, Q, q_I, \Delta, F}$.
    %
    The proof is concluded with the following two claims.
    %
    \begin{claim*}
        For every $1 \leq i \leq k$ and $w \in \Sigma^*$, we have
        \begin{align*}
            g_i(w) = \pi_i(q_I \cdot w^R).
        \end{align*}
    \end{claim*}
    \begin{proof}
        We proceed by induction on $w$.
        In the base case $w = \e$, by definition we have
        $\pi_i(q_I \cdot \e^R) = \pi_i(q_I) = g_i(\e)$.
        In the inductive case, we have
        %
        \begin{align*}
            g_i(a \cdot w)
                &=\coefficient w (\deriveleft a g_i)
                = \coefficient w (\Delta^a_i(g_1, \dots, g_k))
                = \Delta^a_i(g_1(w), \dots, g_k(w))
                = \Delta^a_i(\pi_1(q_I \cdot w^R), \dots, \pi_k(q_I \cdot w^R))
                = \Delta^a_i(q_I \cdot w^R) = \\
                &= \pi_i((q_I \cdot w^R) \cdot a)
                = \pi_i(q_I \cdot (w^R \cdot a))
                = \pi_i (q_I \cdot (a \cdot w)^R).
            \qedhere
        \end{align*}
    \end{proof}
    \begin{claim*}
        We have $g = \sem {q_I}^R$.
    \end{claim*}
    \begin{proof}
        Consider an arbitrary $w \in \Sigma^*$,
        and we need to show $g(w) = F(q_I \cdot w^R)$.
        %
        Thanks to the previous claim, we have
        %
        $g(w)
                = F(g_1(w), \dots, g_k(w))
                = F(\pi_1(q_I \cdot w^R), \dots, \pi_k(q_I \cdot w^R))
                = F(q_I \cdot w^R)$.
    \end{proof}
    \qedhere
\end{proof}
% lax end

% \section{stuff}

% A \emph{series bisimulation} is a binary relation ${\rel} \subseteq \series \Q \Sigma \times \series \Q \Sigma$
% \st~whenever $f \rel g$ we have
% %
% \begin{enumerate}
%     \item $\coefficient \e f = \coefficient \e g$, and
%     \item $\deriveleft a f \rel \deriveleft a g$, for every $a \in \Sigma$.
% \end{enumerate}
% %
% An induction on the length of finite words shows that $f \rel g$ implies $f = g$.
% % when ${\rel}$ is a bisimulation.
% %
% For a relation ${\rel}$, let ${\tilderel}$ be the smallest congruence containing ${\rel}$.
% In other words, $f_0 \tilderel f_1, g_0 \tilderel g_1$ implies
% %
% \begin{enumerate}
%     \item $(c \cdot f_0) \tilderel (c \cdot f_1)$, for all $c \in \Q$,
%     \item $(f_0 + g_0) \tilderel (f_1 + g_1)$, and
%     \item $(f_0 \product g_0) \tilderel (f_1 \product g_1)$
% \end{enumerate}
% %
% A \emph{series bisimulation up-to-context} is a binary relation ${\rel} \subseteq \series \Q \Sigma \times \series \Q \Sigma$
% \st~whenever $f \rel g$ we have
% %
% \begin{enumerate}
%     \item $\coefficient \e f = \coefficient \e g$, and
%     \item $\deriveleft a f \tilderel \deriveleft a g$, for every $a \in \Sigma$.
% \end{enumerate}
% %
% A structural induction on polynomial expressions shows that if $\rel$ is a bisimulation up-to-context,
% then $\tilderel$ is a bisimulation.
% %
% Consequently we obtain the following coinductive reasoning principle.
% %
% \begin{lemma}
%     If $\rel$ is a bisimulation up-to-context and $f \rel g$ then $f = g$.
% \end{lemma}